\section{Methods and Experimental Setup}
\label{sec:methods}

\subsection{Problem formulation}
We study learned conditional channel simulation.
Let \(x \in \mathbb{R}^n\) denote the transmitted channel input and \(y \in \mathbb{R}^n\) the corresponding channel output.
Rather than model \(p(y \mid x)\) directly, we work in residual space and define
\begin{equation}
e = y - x.
\end{equation}
All generative models are therefore trained to approximate the conditional residual law \(p(e \mid x)\), and channel outputs are reconstructed as
\begin{equation}
\hat{y} = x + \hat{e}.
\end{equation}
This parameterization keeps the evaluation target aligned across additive, multiplicative, and nonlinear channels, and it is the representation used throughout the benchmark code.

\subsection{Conditional diffusion baseline}
The diffusion baseline follows the channel-learning setting of \cite{paper2309} and uses an epsilon-prediction DDPM objective \cite{ddpm} with DDIM sampling for deterministic inference trajectories \cite{ddim}.
The backbone is a conditional multilayer perceptron defined by
\begin{equation}
\epsilon_\theta(e_t, t, x),
\end{equation}
where \(e_t\) is the noised residual at diffusion step \(t\), and \(x\) is the channel condition.

The network architecture is a fully connected model with three hidden conditional layers and one linear output head.
In the benchmark implementation, the hidden width is \(N=128\), the residual dimension is \(n=2\), and the first layer takes the concatenated vector \([e_t, x] \in \mathbb{R}^{2n}\).
Each hidden layer is modulated by a learned embedding of the diffusion step \(t\), and a softplus nonlinearity is applied after each conditional affine transform.
The final layer maps the hidden representation back to \(\mathbb{R}^n\) and predicts the diffusion noise on the residual target.

Training uses a cosine noise schedule with horizon \(T=100\) in the publication-grade configuration.
Given residual \(e_0\), a random diffusion step \(t\) is sampled, a noised sample \(e_t\) is formed, and the denoiser is trained with mean-squared error on the true noise:
\begin{equation}
\mathcal{L}_{\text{DDPM}}
=
\mathbb{E}_{e_0, x, t, \epsilon}
\left[
\left\|
\epsilon - \epsilon_\theta(e_t, t, x)
\right\|_2^2
\right].
\end{equation}
At inference time, we report both the stochastic DDPM reverse chain and the deterministic DDIM trajectory.
The main comparison in the final figures uses DDPM(100) and DDIM(100); fast DDIM with fewer steps is kept only as an auxiliary ablation because it distorts the plot scale.

\subsection{Conditional drifting baseline}
The drifting baseline is a conditional one-shot generator inspired by \cite{drifting}.
Following the notation of the drifting paper, let \(f_\theta\) map \(z \sim p_z\) to a sample \(u = f_\theta(z)\), and let
\begin{equation}
q_\theta = (f_\theta)_\# p_z
\end{equation}
denote the induced pushforward distribution.
The central paper-level object is a drifting field \(V_{p,q}(u)\) that indicates how samples from the current model distribution \(q\) should move toward the target distribution \(p\).
The intended equilibrium property is that the field vanishes when the two distributions match:
\begin{equation}
q = p
\quad \Longrightarrow \quad
V_{p,q}(u) = 0.
\end{equation}
This viewpoint leads to a fixed-point style training update.
Starting from \(u = f_\theta(z)\), a drifted target is formed as
\begin{equation}
\tilde{u}
=
\mathrm{stopgrad}\!\left(
u + \eta V_{p,q_\theta}(u)
\right),
\end{equation}
and the generator is optimized by
\begin{equation}
\mathcal{L}_{\text{drift}}
=
\mathbb{E}_{z \sim p_z}
\left[
\left\|
f_\theta(z) - \tilde{u}
\right\|_2^2
\right].
\end{equation}
This makes clear why inference is naturally one-shot: the iterative correction is absorbed into training, and sampling after training requires only one evaluation of \(f_\theta\) \cite{drifting,driftingproj}.

In the channel-learning benchmark, the sample variable above corresponds to the residual \(e = y - x\), and conditioning enters through the transmitted symbol \(x\).
The implemented generator is therefore
\begin{equation}
g_\theta(x, z),
\qquad
z \sim \mathcal{N}(0, I),
\end{equation}
and predicts the residual \(\hat{e} = g_\theta(x, z)\).
In the current implementation, \(g_\theta\) is a two-hidden-layer multilayer perceptron with hidden width \(128\), latent dimension \(16\), and SiLU activations:
\begin{equation}
[x, z]
\rightarrow
\mathrm{Linear}
\rightarrow
\mathrm{SiLU}
\rightarrow
\mathrm{Linear}
\rightarrow
\mathrm{SiLU}
\rightarrow
\mathrm{Linear}.
\end{equation}

The training target is constructed from an attraction--repulsion drift field in residual space.
This is our concrete channel-specific instantiation of the generic field \(V_{p,q}\) from the drifting formulation.
Given generated samples \(G = \{g_i\}\) and positive samples \(P = \{p_j\}\), we compute
\begin{equation}
V(g_i)
=
\left(
\sum_j w^{(P)}_{ij} p_j - g_i
\right)
-
\lambda
\left(
\sum_k w^{(G)}_{ik} g_k - g_i
\right),
\end{equation}
where \(w^{(P)}\) and \(w^{(G)}\) are row-normalized RBF kernel weights.
The first term attracts generated samples toward nearby positive samples, while the second term repels samples away from local generator collapse.
The training target is detached:
\begin{equation}
\tilde{g}_i
=
\mathrm{stopgrad}\!\left(g_i + \eta V(g_i)\right),
\end{equation}
and the generator is trained by
\begin{equation}
\mathcal{L}_{\text{drift}}
=
\frac{1}{N}
\sum_i
\left\|
g_i - \tilde{g}_i
\right\|_2^2.
\end{equation}
Thus, the generic paper-level update \(u \mapsto u + \eta V_{p,q}(u)\) becomes a kernel-based residual update estimated directly from minibatches of generated and true channel residuals.

The corrected benchmark uses the following hyperparameters:
\(\eta = 1.0\),
repulsive weight \(\lambda = 1.0\),
minimum kernel bandwidth \(0.2\),
and maximum drift norm \(2.0\).
Unlike diffusion, inference requires only one forward pass through \(g_\theta(x, z)\), so the per-sample inference cost does not scale with \(T\).

\subsection{Conditional GAN baseline}
The GAN baseline is a conditional residual generator adapted from the \texttt{DM\_OptFib} code path.
The generator receives the concatenated vector \([x, z]\), where \(z \sim \mathcal{N}(0, I)\), and outputs the full channel output \(\hat{y}\).
Architecturally, the generator uses an input linear layer, two batch-normalized hidden blocks, a central hidden layer, two additional batch-normalized hidden blocks, and a residual skip from the condition input \(x\) to the output.
The discriminator takes the concatenated pair \([y, x]\), processes it with a residual-style multilayer perceptron, concatenates the original input back into the hidden state, and produces a scalar realism score.

Two adversarial objectives are implemented in the benchmark code: WGAN-GP and a stronger feature-aligned GAN mode denoted \texttt{GAN\_FA}.
The final benchmark uses \texttt{GAN\_FA}, since it was substantially more stable than the plain WGAN-GP baseline in this channel-learning setting.
In \texttt{GAN\_FA}, the discriminator is trained with BCE-with-logits and label smoothing, and the generator is trained with an adversarial term plus an \(L_1\) reconstruction term:
\begin{equation}
\mathcal{L}_{G}
=
\mathcal{L}_{\text{adv}}
+
\alpha
\mathbb{E}\left[\|\hat{y} - y\|_1\right].
\end{equation}
In the current implementation, \(\alpha = 1.0\), generator updates are delayed relative to discriminator updates, and the default hidden width is \(128\).

\subsection{Channel models}
We benchmark four channels.

\textbf{AWGN:}
\begin{equation}
y = x + n,
\qquad
n \sim \mathcal{N}(0, \sigma^2 I).
\end{equation}

\textbf{Rayleigh:}
the transmitted symbol is multiplied elementwise by a Rayleigh fading coefficient and then corrupted by additive Gaussian noise.

\textbf{SSPA:}
the input amplitude is passed through a smooth solid-state power amplifier nonlinearity before additive noise is applied.

\textbf{OptFib:}
we include a simplified memoryless optical-fiber proxy channel adapted from \texttt{DM\_OptFib}.
The I/Q pair is propagated through \(K_{\text{step}}\) repeated nonlinear phase rotations with per-step Gaussian noise injection.
In the current benchmark, the default parameters are \(L = 5000\), \(\gamma = 1.27\), \(K_{\text{step}} = 50\), and \(P_n = -21.3\) dBm.

\subsection{Training protocol}
The publication runner is implemented in \texttt{publication\_benchmark\_t100.py} and is designed for multi-seed GPU experiments.
The publication-scale default configuration is:
\begin{itemize}
\item diffusion horizon \(T=100\),
\item residual dimension \(n=2\),
\item noise standard deviation \(0.3\),
\item training set size \(120{,}000\),
\item diffusion epochs \(60\),
\item GAN epochs \(120\),
\item batch size \(512\),
\item evaluation set size \(20{,}000\),
\item GAN mode \texttt{GAN\_FA}.
\end{itemize}
The drifting model uses the same channel samples and residual-space target as the diffusion and GAN baselines.
To support publication-ready reporting, the runner performs multi-seed sweeps, stores per-seed outputs, computes mean and standard deviation, estimates a \(95\%\) bootstrap confidence interval for SWD, and records pairwise method wins across seeds.

\subsection{Evaluation metric}
The primary metric is sliced Wasserstein distance (SWD) on residual vectors:
\begin{equation}
\mathrm{SWD}\!\left(\{\hat{e}_i\}, \{e_i\}\right).
\end{equation}
In the local benchmark code, SWD is estimated with \(256\) random one-dimensional projections.
Residual-space SWD was selected over histogram \(L_1\) distance because it is better behaved across channel families and captures geometry in \(\mathbb{R}^2\) rather than only marginal overlap.

\subsection{Inference-time comparison}
The latency comparison is measured separately from the fidelity benchmark.
For the stored CPU timing run in this repository, DDPM(100) required \(4.2612\) seconds per repeat over 20,480 samples, while the one-shot drifting generator required \(0.0163\) seconds under the same batch structure.
This corresponds to an observed speedup of approximately \(260.8\times\) in favor of drifting.

\subsection{Reproducibility}
The benchmark script exposes an explicit random seed argument and the publication runner defaults to deterministic execution unless disabled.
Additional sanity checks performed in this repository include: metric identity and symmetry checks for SWD, byte-identical reruns under fixed seed, finite-value validation of stored JSON result files, and small cross-seed ranking checks on AWGN and OptFib.
These checks are not a substitute for the full publication-scale run, but they provide confidence that the current implementation is internally consistent.
